\documentclass[11pt]{article}
\usepackage{ap-physc}
\course{AP Physics C: Mechanics}
\headcourse{AP Physics C: Mechanics}
\headpart{Unit 4 \textperiodcentered\ Question Set}
\title{Unit 4: Linear Momentum: Question Set}
\author{Paul Wang}
\date{November 2026}
\docpart{Unit 4 \textperiodcentered\ Question Set}

\begin{document}
\maketitle

\begin{abstract}
Exam-style practice on Unit 4 (topics 4.1--4.4), which is 10--20\% of the multiple-choice section. Section~1--3 hold 12 four-choice multiple-choice questions grouped by topic; most need a short calculation or a line of reasoning, as on the exam. The last page before the key is a full Experimental Design and Analysis free-response question (10 points). All items are original and written for practice. Use the AP equation sheet and a calculator. Allow about 24 minutes for the multiple choice (the exam's pace is about 2 minutes per question) and 24 minutes for the free response. The answer key explains every item and names its topic; a scoring guide for the free response follows it.
\end{abstract}

\scoringline{MCQ: \blank\ / \qtotal \qquad FRQ: \blank\ / 10 \qquad Target: 8+ MCQ \qquad Pace: 2 min per MCQ}

\section{Impulse (4.1--4.2)}

\begin{q}{2}{4.2 Impulse of a variable force}{$J = \int_0^2 30t^2\,dt = 10t^3\big|_0^2 = 80$ N$\cdot$s, so $v = J/m = 80/3.0 = 27$ m/s.}
A net force $F = 30t^2$ (newtons, $t$ in seconds) acts for 2.0 s on a 3.0 kg cart that starts from rest. What is the cart's final speed?
\chh{$20$ m/s}{$27$ m/s}{$40$ m/s}{$80$ m/s}
\end{q}
\begin{q}{4}{4.2 Average force}{$|\Delta p| = 0.15(40 + 50) = 13.5$ kg$\cdot$m/s, so $\bar F = 13.5/0.0050 = 2700$ N.}
A 0.15 kg baseball arrives at 40 m/s and leaves the bat at 50 m/s in the opposite direction. The contact lasts 0.0050 s. What is the average force of the bat on the ball?
\chh{$300$ N}{$1200$ N}{$1500$ N}{$2700$ N}
\end{q}
\begin{q}{1}{4.2 Variable mass}{Each second 2.0 kg of sand must be brought from zero to 3.0 m/s horizontally: $F = v\,dm/dt = (3.0)(2.0) = 6.0$ N.}
Sand falls vertically onto a horizontal conveyor belt at 2.0 kg/s. The belt moves at a constant 3.0 m/s. What extra horizontal force must the motor supply to keep the belt's speed constant? (Ignore friction in the rollers.)
\chh{$6.0$ N}{$3.0$ N}{$1.5$ N}{$0$ N}
\end{q}
\begin{q}{3}{4.2 Impulse and contact time}{The egg's momentum goes from $mv$ to 0 either way, so the impulse is the same. Spreading it over a longer time lowers $\bar F = \Delta p/\Delta t$.}
An egg dropped onto a pillow survives, while the same egg dropped from the same height onto concrete breaks. Compared with the concrete, the pillow
\ch{gives a smaller impulse over a longer time}{gives a smaller impulse over the same time}{gives the same impulse over a longer time, so the average force is smaller}{gives the same impulse but a larger average force}
\end{q}

\section{Conservation of momentum (4.3)}

\begin{q}{2}{4.3 Recoil}{Momentum stays zero: $60\,v = 2.0(15)$, so $v = 0.50$ m/s backward.}
A 60 kg skater at rest on frictionless ice throws a 2.0 kg ball forward at 15 m/s. How fast does the skater recoil?
\chh{$0.033$ m/s}{$0.50$ m/s}{$2.7$ m/s}{$15$ m/s}
\end{q}
\begin{q}{3}{4.3 Center-of-mass velocity}{$v_{cm} = \dfrac{\sum m_iv_i}{\sum m_i} = \dfrac{(2.0)(3.0) + (1.0)(-3.0)}{3.0} = 1.0$ m/s east, before and after any collision.}
A 2.0 kg cart moves east at 3.0 m/s toward a 1.0 kg cart moving west at 3.0 m/s. What is the velocity of the two-cart system's center of mass?
\chh{$3.0$ m/s east}{$2.0$ m/s east}{$1.0$ m/s east}{$0$ m/s}
\end{q}
\begin{q}{4}{4.3 Explosion in flight}{The explosion's forces are internal, so the center of mass follows the original parabola. The second piece lands at 1.5 times the original range, keeping the center of mass on track.}
A shell fired from level ground explodes at the top of its arc into two pieces of equal mass. One piece falls straight down from rest. Where does the center of mass of the two pieces land (assuming both land at the same time)?
\ch{Directly below the explosion}{Halfway between the launch point and the explosion point}{Twice as far as the shell would have gone}{Where the shell would have landed without exploding}
\end{q}

\section{Collisions (4.4)}

\begin{q}{1}{4.4 Perfectly inelastic}{Momentum: $1500(20) = 4000\,v$, so $v = 7.5$ m/s.}
A 1500 kg car moving at 20 m/s rear-ends a 2500 kg truck at rest, and the two lock together. What is their speed just after the collision?
\chh{$7.5$ m/s}{$10$ m/s}{$12$ m/s}{$20$ m/s}
\end{q}
\begin{q}{4}{4.4 Energy in a perfectly inelastic collision}{For a target at rest, $K_f/K_i = m_1/(m_1 + m_2) = 1500/4000 = 0.375$, so 62.5\% is lost.}
In the collision of the previous question, what percentage of the car's kinetic energy is lost?
\chh{$0\%$}{$38\%$}{$50\%$}{$63\%$}
\end{q}
\begin{q}{3}{4.4 Elastic collision}{$v_{2f} = \dfrac{2m_1}{m_1 + m_2}v_{1i} = \dfrac{4.0}{6.0}(6.0) = 4.0$ m/s; the 2.0 kg ball bounces back at 2.0 m/s.}
A 2.0 kg ball moving at 6.0 m/s collides elastically head-on with a 4.0 kg ball at rest. What is the 4.0 kg ball's speed afterward?
\chh{$2.0$ m/s}{$3.0$ m/s}{$4.0$ m/s}{$6.0$ m/s}
\end{q}
\begin{q}{2}{4.4 Two-dimensional collision}{Conserve each component. A afterward: $(1.8, 2.4)$ m/s. B must carry $(5.0 - 1.8, -2.4) = (3.2, -2.4)$ m/s, a speed of 4.0 m/s.}
Puck A (0.50 kg) slides east at 5.0 m/s and strikes identical puck B at rest. Afterward A moves at 3.0 m/s, $53^\circ$ north of east. What is B's speed afterward?
\chh{$8.0$ m/s}{$4.0$ m/s}{$3.2$ m/s}{$2.0$ m/s}
\end{q}
\begin{q}{1}{4.4 Ballistic pendulum}{Swing (energy): $V = \sqrt{2gh} = 1.98$ m/s. Collision (momentum): $v = \dfrac{(m + M)V}{m} = \dfrac{2.01(1.98)}{0.010} = 400$ m/s.}
A 10 g bullet embeds in a 2.0 kg block hanging from strings. The block and bullet swing up 0.20 m. What was the bullet's speed? ($g = 9.8\ \mathrm{m/s^2}$)
\chh{$400$ m/s}{$200$ m/s}{$28$ m/s}{$2.0$ m/s}
\end{q}

\section{Free response}

\begin{frq}{Question 1 \textperiodcentered\ Experimental Design and Analysis \textperiodcentered\ 10 points \textperiodcentered\ 24 minutes}
Students study a collision on a level, low-friction track. Cart A, of mass $m = 0.50$ kg, is pushed so that it moves at an unknown speed $v_0$ toward cart B, of mass $M$, which is at rest. Velcro makes the carts stick together. The students repeat the trial with different values of $M$, launching cart A the same way each time, and record the speed $v_f$ of the joined carts.

\medskip
\begin{tabular}{@{}lccccc@{}}\toprule
$M$ (kg) & 0.25 & 0.50 & 0.75 & 1.00 & 1.25 \\ \midrule
$v_f$ (m/s) & 0.54 & 0.41 & 0.31 & 0.27 & 0.23 \\ \bottomrule
\end{tabular}

\begin{enumerate}[label=(\alph*),leftmargin=2em,itemsep=4pt]
\item Describe a procedure the students could use to measure $v_f$, including the equipment and how each measurement is used. Identify one source of experimental uncertainty and its effect on the measured $v_f$. \pts{3}
\item Indicate quantities the students could graph to produce a straight line, and explain how that line relates to $v_0$. Use conservation of momentum. \pts{2}
\item Using the data, calculate the values to be graphed. Determine the slope of a best-fit line (by plotting or by a calculator fit). \pts{3}
\item Use the slope to calculate an experimental value of $v_0$. Then state whether the intercept of your line is consistent with that value. \pts{2}
\end{enumerate}
\end{frq}

\makekey

\subsection*{Free-response scoring guide (10 points)}
\small
Award each point only when the answer does what the model does. Points are independent unless a part says otherwise; a correct method with an answer carried forward from an earlier error earns the later points.
\begin{scoring}
(a) (3) & 1 pt: a valid speed measurement after the collision, e.g.\ a motion sensor or video with a ruler in frame, or a photogate timing a flag of known length. 1 pt: speed obtained correctly, e.g.\ slope of the position--time data after the collision, or flag length $\div$ gate time. 1 pt: one real source with its direction, e.g.\ friction slows the carts, so the measured $v_f$ is lower than the true speed just after impact. & Listing ``human error'' or ``measure carefully'' with no mechanism or direction. \\
(b) (2) & 1 pt: $mv_0 = (m + M)v_f$, so $\dfrac{1}{v_f} = \dfrac{1}{v_0} + \dfrac{1}{mv_0}M$. 1 pt: graph $1/v_f$ versus $M$: slope $= 1/(mv_0)$ (intercept $1/v_0$). & Graphing $v_f$ versus $M$, which is a curve. \\
(c) (3) & 1 pt: $1/v_f = 1.85,\ 2.44,\ 3.23,\ 3.70,\ 4.35$ s/m. 1 pt: a reasonable straight best-fit line (on a plot) or a least-squares fit. 1 pt: slope $\approx 2.5$ s/(m$\cdot$kg), accept $2.3$--$2.7$. & Calculating the slope from two data points instead of the line. \\
(d) (2) & 1 pt: $v_0 = 1/(m \times \text{slope}) = 1/(0.50 \times 2.5) = 0.80$ m/s. 1 pt: intercept $\approx 1.24$ s/m, so $1/\text{intercept} \approx 0.81$ m/s, consistent with $0.80$ m/s. & Forgetting the factor $m$ in the slope. \\
\end{scoring}

\end{document}
